\textbf{Séries de Taylor (dois termos) para $\sinh x$ e $\cosh x$:}

Das definições exponenciais:
\[
\sinh x = \frac{e^x-e^{-x}}{2},\qquad \cosh x = \frac{e^x+e^{-x}}{2}.
\]

Usando $e^x=1+x+x^2/2!+x^3/3!+\cdots$:
\[
e^x-e^{-x} = 2x + \frac{2x^3}{3!}+\frac{2x^5}{5!}+\cdots \implies \sinh x = x+\frac{x^3}{6}+\frac{x^5}{120}+\cdots\approx x+\frac{x^3}{6},
\]
\[
e^x+e^{-x} = 2+\frac{2x^2}{2!}+\frac{2x^4}{4!}+\cdots \implies \cosh x = 1+\frac{x^2}{2}+\frac{x^4}{24}+\cdots\approx 1+\frac{x^2}{2}.
\]

\textbf{Computação de $1-\sinh x$ e $1-\cosh x$:}
\[
1-\sinh x = 1 - x - \frac{x^3}{6} + \cdots,
\]
\[
1-\cosh x = -\frac{x^2}{2} - \frac{x^4}{24} + \cdots \approx -\frac{x^2}{2}.
\]

\textbf{Comparação com as funções trigonométricas:}

\begin{center}
\begin{tabular}{ll}
\hline
Circular & Hiperbólica \\
\hline
$\sin x \approx x - x^3/6$ & $\sinh x \approx x + x^3/6$ \\
$\cos x \approx 1 - x^2/2$ & $\cosh x \approx 1 + x^2/2$ \\
$1-\sin x \approx 1-x+x^3/6$ & $1-\sinh x \approx 1-x-x^3/6$ \\
$1-\cos x \approx x^2/2$ & $1-\cosh x \approx -x^2/2$ \\
\hline
\end{tabular}
\end{center}

Note que o sinal de $x^2/2$ se inverte: enquanto $\cos x<1$ para $x\ne0$, temos $\cosh x>1$ para todo $x\ne0$.